% © 2026 Ing. David Jaroš — CC BY-NC-ND 4.0
%
% NCG Spectral Action → B₀ Derivation
% Výsledky přímého výpočtu Poisson resummation + KK redukce
% Klasifikace: [L1] pro KK argument, [MC] pro N_eff identifikaci

\documentclass[12pt,a4paper]{article}
\usepackage{amsmath,amssymb,amsthm}
\usepackage{booktabs}
\usepackage{tcolorbox}
\tcbuselibrary{skins,breakable}
\newtheorem{theorem}{Theorem}[section]
\newtheorem{lemma}[theorem]{Lemma}
\newtheorem{proposition}[theorem]{Proposition}
\theoremstyle{remark}
\newtheorem{remark}[theorem]{Remark}
\newcommand{\PROVED}{\textbf{[Proved]}}
\newcommand{\MC}{\textbf{[MC]}}
\newcommand{\OPEN}{\textbf{[Open]}}
\newcommand{\Lone}{\textbf{[L1]}}
\newcommand{\Lzero}{\textbf{[L0]}}
\newcommand{\STD}{\textbf{[STD]}}

\title{\textbf{NCG Spectral Action and the Origin of $B_0$}\\
{\large Kaluza--Klein Derivation of $B_0 = \tfrac{2\pi}{3}N_{\mathrm{eff}}$}}
\author{Ing.\ David Jaro\v{s}}
\date{2026-05-11}

\begin{document}
\maketitle

\begin{abstract}
We derive $B_0 = \tfrac{2\pi}{3}N_{\mathrm{eff}}$ from the NCG spectral action
on $M^4\times S^1_\psi$ via Kaluza--Klein reduction and heat-kernel asymptotics.
The logarithmic term $B_0\ln\Lambda$ in the effective action arises from 4D
one-loop diagrams for each winding mode $\Theta_n$; summation over modes yields
the $N_{\mathrm{eff}}$ prefactor.
We also identify the UBT winding partition function
$Z_{\mathrm{wind}}(\tau) = \vartheta_3(0\mid i\tau)$ as the natural
spectral-action object, and note that its modular transformation
generates a $\ln\Lambda$ term whose coefficient is a structural
bridge candidate to $B_0$.  Gap~G137-B remains open.
\end{abstract}

%──────────────────────────────────────────────────────────────────────────────
\section{Setup: Spectral Triple on $M^4\times S^1_\psi$}
\label{sec:setup}
%──────────────────────────────────────────────────────────────────────────────

The UBT Dirac operator on $M^4\times S^1_\psi$ decomposes as
\begin{equation}
  D_{M^4\times S^1}^2 = D_{M^4}^2 + (-i\partial_\psi)^2.
  \label{eq:D2}
\end{equation}
The spectrum of $(-i\partial_\psi)^2$ on $S^1_\psi$ consists of
Kaluza--Klein masses
\begin{equation}
  m_n^2 = \frac{n^2}{R_\psi^2}, \qquad n\in\mathbb{Z}.
  \label{eq:KKspectrum}
\end{equation}
The spectral action is therefore
\begin{equation}
  \mathrm{Tr}\!\left[f\!\left(\frac{D}{{\Lambda}}\right)\right]
  = \sum_{n\in\mathbb{Z}}
    \mathrm{Tr}_{4D}\!\left[f\!\left(\frac{\sqrt{D_{M^4}^2 + m_n^2}}{\Lambda}\right)\right].
  \label{eq:spectral_KK}
\end{equation}

%──────────────────────────────────────────────────────────────────────────────
\section{Logarithmic Term via Heat-Kernel Asymptotics \STD}
\label{sec:heatkernel}
%──────────────────────────────────────────────────────────────────────────────

For each KK mode $n$, the 4D spectral action with a massive field of
mass $m_n$ has the Seeley--DeWitt asymptotic expansion
\begin{equation}
  \mathrm{Tr}_{4D}\!\left[f\!\left(\frac{\sqrt{D_{M^4}^2+m_n^2}}{\Lambda}\right)\right]
  \sim \Lambda^4 f_4 + \Lambda^2 f_2\,m_n^2
    + \underbrace{\beta_{\mathrm{1-loop}}\cdot\ln\Lambda}_{\text{log term}}
    + O(\Lambda^{-2}),
  \label{eq:SdW}
\end{equation}
where the logarithmic coefficient for a complex scalar in 4D is the
one-loop $\beta$-function contribution
\begin{equation}
  \beta_{\mathrm{1-loop}} = \frac{e^2}{24\pi^2}
  \qquad \text{(independent of }m_n\text{ for }\Lambda\gg m_n).
  \label{eq:beta1loop}
\end{equation}
The key observation is that $\beta_{\mathrm{1-loop}}$ is \emph{independent
of the mass} $m_n$ in the UV regime $\Lambda\gg m_n$.

\begin{proposition}[Logarithmic term from KK sum \STD]
\label{prop:log}
The total $\ln\Lambda$ coefficient in the UBT spectral action is
\begin{equation}
  \boxed{
  \frac{\partial}{\partial\ln\Lambda}\,
  \mathrm{Tr}\!\left[f\!\left(\frac{D}{\Lambda}\right)\right]
  = N_{\mathrm{eff}}(\Lambda)\cdot\beta_{\mathrm{1-loop}}
  = \frac{e^2}{24\pi^2}\cdot N_{\mathrm{eff}}(\Lambda),}
  \label{eq:logcoeff}
\end{equation}
where $N_{\mathrm{eff}}(\Lambda) = \#\{n\in\mathbb{Z}\mid m_n < \Lambda\}
= 2\lfloor R_\psi\Lambda\rfloor + 1$ is the number of KK modes below the
UV cutoff.
\end{proposition}

\begin{proof}
Sum~\eqref{eq:spectral_KK} over modes $n$ with $m_n<\Lambda$; each
contributes $\beta_{\mathrm{1-loop}}\cdot\ln\Lambda$ by \eqref{eq:SdW}--\eqref{eq:beta1loop}.
Modes with $m_n>\Lambda$ are Boltzmann-suppressed by $e^{-m_n^2/\Lambda^2}\to 0$.
\end{proof}

\begin{remark}[Origin of $B_0$]
In the UBT convention $\partial(1/\alpha)/\partial\ln\Lambda = B_0/(2\pi)$,
Proposition~\ref{prop:log} gives
\begin{equation}
  B_0 = \frac{2\pi}{3}\,N_{\mathrm{eff}}.
  \label{eq:B0NCG}
\end{equation}
This is precisely the relation derived by independent methods in
\texttt{canonical/n\_eff/step2\_vacuum\_polarization.tex} (\textbf{[L1]}).
The NCG derivation is an independent confirmation at the same level.
\end{remark}

%──────────────────────────────────────────────────────────────────────────────
\section{Winding Partition Function and Theta Functions \MC}
\label{sec:theta}
%──────────────────────────────────────────────────────────────────────────────

The winding-mode partition function on $S^1_\psi$ is
\begin{equation}
  Z_{\mathrm{wind}}(\tau)
  = \sum_{n\in\mathbb{Z}} e^{-\pi\tau n^2}
  = \vartheta_3(0\mid i\tau),
  \qquad \tau = \frac{1}{\pi R_\psi^2\Lambda^2}.
  \label{eq:Zwind}
\end{equation}
The logarithmic effective action is
$S_{\mathrm{log}} = \ln Z_{\mathrm{wind}}(\tau)$.
Under the modular transformation $\tau\to 1/\tau$:
\begin{equation}
  \vartheta_3(0\mid i/\tau) = \sqrt{\tau}\,\vartheta_3(0\mid i\tau),
\end{equation}
so for small $\tau$ (large $\Lambda$):
\begin{equation}
  Z_{\mathrm{wind}} \sim \tau^{-1/2}
  = \sqrt{\pi}\,R_\psi\Lambda
  \qquad\Longrightarrow\qquad
  S_{\mathrm{log}} \sim \ln\Lambda.
  \label{eq:logfromtheta}
\end{equation}

\begin{remark}[Structural bridge candidate \MC]
Equation~\eqref{eq:Zwind} identifies $Z_{\mathrm{wind}}$ with the
Jacobi theta function $\vartheta_3$, the same object that appears in the
candidate formula $B = N_{\mathrm{eff}}^{3/2}(2\eta(i))^{1/4}$ studied
in \texttt{research\_tracks/T3\_ALPHA/eta\_i\_alpha\_coefficient\_derivation.tex}.
Specifically, $\eta(\tau)^3 = \vartheta_2\vartheta_3\vartheta_4/(2\cdot 4)$
at $\tau=i$.  This algebraic overlap is a structural coincidence
classified \MC\ pending a derivation linking the spectral-action
normalization to $B_{\mathrm{phenom}}$.
\end{remark}

%──────────────────────────────────────────────────────────────────────────────
\section{Vacuum Polarization from Background Field Method \STD+MC}
\label{sec:bgfield}
%──────────────────────────────────────────────────────────────────────────────

With EM background $A_\mu$, the winding spectrum shifts to
$\lambda_n + eA\cdot n$, so
\begin{equation}
  Z_{\mathrm{wind}}(\tau, z) = \vartheta_3(z\mid i\tau),
  \quad z = \frac{eAR_\psi}{\pi}.
\end{equation}
The vacuum polarization is
\begin{equation}
  \Pi \propto e^2 R_\psi^2\,\frac{\vartheta_3''(0\mid i\tau)}{\vartheta_3(0\mid i\tau)}.
  \label{eq:vacpol}
\end{equation}
Numerically, for small $\tau$:
\begin{equation}
  \frac{\vartheta_3''(0\mid i\tau)}{\vartheta_3(0\mid i\tau)}
  \approx -\frac{2\pi^2}{\tau}
  \qquad[\text{verified numerically}].
\end{equation}
With $\tau = 1/(\pi R_\psi^2\Lambda^2)$ this gives $\Pi\sim\Lambda^2$
(quadratic divergence), not the logarithmic $B_0\ln\Lambda$.
The logarithmic term requires the full 4D loop integral and
comes from the $a_4$ Seeley--DeWitt coefficient, confirming the
KK result of Section~\ref{sec:heatkernel}.

%──────────────────────────────────────────────────────────────────────────────
\section{Summary and Gap Status}
\label{sec:summary}
%──────────────────────────────────────────────────────────────────────────────

\begin{tcolorbox}[colback=green!4!white,colframe=green!55!black,
  title={NCG derivation of $B_0$ — status summary}]
\begin{center}
\begin{tabular}{lll}
\toprule
\textbf{Result} & \textbf{Level} & \textbf{Method} \\
\midrule
$B_0 = \tfrac{2\pi}{3}N_{\mathrm{eff}}$ from KK spectral action
  & \Lone/\STD & Prop.~\ref{prop:log} \\
$\ln\Lambda$ from 4D loops, not Poisson sum
  & \Lone/\STD & heat kernel \\
$Z_{\mathrm{wind}} = \vartheta_3(0\mid i\tau)$ identification
  & \Lone & Eq.~\eqref{eq:Zwind} \\
Modular $\to\ln\Lambda$ via $\vartheta_3$ transformation
  & \MC & Eq.~\eqref{eq:logfromtheta} \\
$\vartheta_3$ links to $\eta(i)$ formula for $B$
  & \MC\ (structural) & Remark in §\ref{sec:theta} \\
\midrule
$B_0\to B_{\mathrm{phenom}}$ bridge (Gap G137-B)
  & \OPEN & not derived \\
$N_{\mathrm{eff}}=12$ from NCG (vs.\ audit $N_{\mathrm{eff}}=3$)
  & \OPEN & not derived \\
\bottomrule
\end{tabular}
\end{center}
\end{tcolorbox}

\paragraph{What is new.}
The NCG KK derivation of $B_0=\tfrac{2\pi}{3}N_{\mathrm{eff}}$ is an
independent confirmation at \Lone/\STD\ level, consistent with
\texttt{step2\_vacuum\_polarization.tex}.
The identification $Z_{\mathrm{wind}}=\vartheta_3$ is a new structural
observation linking the winding partition function to the theta-function
machinery of the alpha track.

\paragraph{What is not new.}
Gap G137-B (derive $B_{\mathrm{phenom}}\approx 46.3$) and the $N_{\mathrm{eff}}$
audit gap remain open.  The NCG framework clarifies their origin but does
not resolve them.

\end{document}
